\setcounter{framenumber}{0}
%27

\begin{frame}
  \titlepage
\end{frame}

\begin{frame}{Learning Goals}
By the end of this lecture, students should be able to:
\begin{itemize}
    \item define the Beta distribution;
    \item understand how Beta parameters affect shape;
    \item recognize Uniform distributions as special Beta distributions;
    \item define the Gamma function and Gamma distribution;
    \item interpret Gamma distributions as sums of Exponential random variables;
    \item compute means and variances of Beta and Gamma random variables;
    \item understand the relationship between Gamma and Poisson processes.
\end{itemize}
\end{frame}

\lecturepart{Beta Distribution}

\begin{frame}{Why the Beta Distribution?}
Many random quantities naturally live between
\[
0
\quad \text{and} \quad
1.
\]

Examples:
\begin{itemize}
    \item probabilities;
    \item proportions;
    \item percentages;
    \item fractions;
    \item success rates.
\end{itemize}

\vspace{0.25cm}

\begin{keybox}{Main idea}
The Beta distribution is one of the most important distributions on the interval
\[
(0,1).
\]
\end{keybox}
\end{frame}

\begin{frame}{Beta Function}
\begin{keybox}{Definition}
For
\[
a,b>0,
\]
the Beta function is
\[
{
B(a,b)
=
\int_0^1
x^{a-1}(1-x)^{b-1}\,dx.
}
\]
\end{keybox}

\vspace{0.25cm}

The Beta function serves as a normalization constant.
\end{frame}

\begin{frame}{Definition of the Beta Distribution}
\begin{keybox}{Definition}
A random variable $X$ has the Beta distribution with parameters
\[
a,b>0
\]
if its PDF is
\[
{
f(x)
=
\frac{
x^{a-1}(1-x)^{b-1}
}{
B(a,b)
},
\qquad 0<x<1.
}
\]

Notation:
\[
X\sim \operatorname{Beta}(a,b).
\]
\end{keybox}
\end{frame}

\begin{frame}{Shape of Beta Distributions}
Different parameter values create very different shapes.

\vspace{0.25cm}

Examples:
\begin{itemize}
    \item
    \[
    a=b=1
    \]
    gives a Uniform distribution;

    \item
    \[
    a,b>1
    \]
    produces mound-shaped densities;

    \item
    \[
    a<1
    \]
    or
    \[
    b<1
    \]
    can create spikes near the endpoints.
\end{itemize}

\vspace{0.25cm}

\begin{keybox}{Important}
The Beta family is extremely flexible.
\end{keybox}
\end{frame}

\begin{frame}{Uniform Distribution as a Beta Distribution}
If
\[
a=b=1,
\]
then
\[
B(1,1)=1,
\]
and the PDF becomes
\[
f(x)=1,
\qquad 0<x<1.
\]

Thus:
\[
{
\operatorname{Beta}(1,1)
=
\operatorname{Unif}(0,1).
}
\]
\end{frame}

\begin{frame}{Mean and Variance of Beta}
If
\[
X\sim \operatorname{Beta}(a,b),
\]
then
\[
{
\mathbb{E}(X)
=
\frac{a}{a+b}.
}
\]

Also,
\[
{
\operatorname{Var}(X)
=
\frac{ab}{(a+b)^2(a+b+1)}.
}
\]

\vspace{0.25cm}

\begin{keybox}{Interpretation}
The parameters control both:
\begin{itemize}
    \item the center;
    \item the concentration.
\end{itemize}
\end{keybox}
\end{frame}

\begin{frame}{Beta Distribution and Beliefs}
Suppose a coin has unknown success probability
\[
p.
\]

Before observing data, we may model uncertainty about $p$ using a Beta distribution.

\vspace{0.25cm}

Example:
\[
p\sim \operatorname{Beta}(2,2)
\]
represents a prior belief favoring values near
\[
0.5.
\]

\vspace{0.25cm}

\begin{keybox}{Bayesian interpretation}
Beta distributions are natural models for uncertain probabilities.
\end{keybox}
\end{frame}

\begin{frame}{Beta-Binomial Connection}
Suppose:
\begin{itemize}
    \item
    \[
    p\sim \operatorname{Beta}(a,b),
    \]
    and conditional on $p$,

    \item
    \[
    X\mid p\sim \operatorname{Bin}(n,p).
    \]
\end{itemize}

Then the posterior distribution of $p$ given $X=k$ is again Beta:

\[
\boxed{
p\mid X=k
\sim
\operatorname{Beta}(a+k,b+n-k).
}
\]

\vspace{0.25cm}

This is called conjugacy.
\end{frame}

\begin{frame}{Bayes' Billiards}
Imagine tossing a ball uniformly onto a billiard table.

A second sequence of balls lands independently.

\vspace{0.25cm}

We ask:
what fraction land to the left of the first ball?

\vspace{0.25cm}

This naturally leads to Beta distributions.

\vspace{0.25cm}

\begin{keybox}{Historical note}
This is one of the earliest Bayesian probability models.
\end{keybox}
\end{frame}

\lecturepart{Gamma Function}

\begin{frame}{Why the Gamma Function?}
Recall:
\[
n!
=
1\cdot2\cdot3\cdots n.
\]

Can factorials be extended to non-integer values?

\vspace{0.25cm}

\begin{keybox}{Answer}
Yes. The Gamma function generalizes factorials to positive real numbers.
\end{keybox}

\begin{keybox}{Definition}
For
\[
\alpha>0,
\]
the Gamma function is
\[
{
\Gamma(\alpha)
=
\int_0^\infty
x^{\alpha-1}e^{-x}\,dx.
}
\]
\end{keybox}
\end{frame}



\begin{frame}{Connection with Factorials}
\begin{keybox}{Theorem}
For positive integers $n$,
\[
\boxed{
\Gamma(n)=(n-1)!.
}
\]
\end{keybox}

\textbf{Proof}. Show $\Gamma(n)=(n-1)\times \Gamma(n-1)$ using integration by parts, and then solve the recursion with induction.


Examples:
\[
\Gamma(1)=1,
\]
\[
\Gamma(2)=1!,
\]
\[
\Gamma(3)=2!,
\]
etc.

\end{frame}



\lecturepart{Gamma Distribution}

\begin{frame}{Definition of Gamma Distribution}
\begin{keybox}{Definition}
A random variable $X$ has the Gamma distribution with parameters
\[
\alpha,\lambda>0
\]
if its PDF is
\[
{
f(x)
=
\frac{
\lambda^\alpha
}{
\Gamma(\alpha)
}
x^{\alpha-1}e^{-\lambda x},
\qquad x>0.
}
\]

Notation:
\[
X\sim \operatorname{Gamma}(\alpha,\lambda).
\]
\end{keybox}
\end{frame}

\begin{frame}{Exponential Distribution as a Gamma Distribution}
If
\[
\alpha=1,
\]
then
\[
\Gamma(1)=1,
\]
and the PDF becomes
\[
f(x)=\lambda e^{-\lambda x}.
\]

Thus:
\[
{
\operatorname{Gamma}(1,\lambda)
=
\operatorname{Exp}(\lambda).
}
\]
\end{frame}
\begin{frame}{Mean and Variance of Gamma}
If
\[
X\sim \operatorname{Gamma}(\alpha,\lambda),
\]
then
\[
\mathbb{E}(X)=\frac{\alpha}{\lambda}.
\]

Also,
\[
\operatorname{Var}(X)=\frac{\alpha}{\lambda^2}.
\]
\end{frame}



\begin{frame}{Proof: Mean and Variance of Gamma}
The density is
\[
f_X(x)
=
\frac{\lambda^\alpha}{\Gamma(\alpha)}
x^{\alpha-1}e^{-\lambda x},
\qquad x>0.
\]

Then
\[
\mathbb{E}(X)
=
\int_0^\infty x f_X(x)\,dx
=
\frac{\lambda^\alpha}{\Gamma(\alpha)}
\int_0^\infty x^\alpha e^{-\lambda x}\,dx.
\]

Using
\[
\int_0^\infty x^\alpha e^{-\lambda x}\,dx
=
\frac{\Gamma(\alpha+1)}{\lambda^{\alpha+1}},
\]
we get
\[
\mathbb{E}(X)
=
\frac{\lambda^\alpha}{\Gamma(\alpha)}
\cdot
\frac{\Gamma(\alpha+1)}{\lambda^{\alpha+1}}.
\]

Since \(\Gamma(\alpha+1)=\alpha\Gamma(\alpha)\),
\[
\mathbb{E}(X)=\frac{\alpha}{\lambda}.
\]
\end{frame}



\begin{frame}{Proof: Variance of Gamma}
First compute
\[
\mathbb{E}(X^2)
=
\int_0^\infty x^2 f_X(x)\,dx.
\]

Thus
\[
\mathbb{E}(X^2)
=
\frac{\lambda^\alpha}{\Gamma(\alpha)}
\int_0^\infty x^{\alpha+1}e^{-\lambda x}\,dx.
\]

Using
\[
\int_0^\infty x^{\alpha+1}e^{-\lambda x}\,dx
=
\frac{\Gamma(\alpha+2)}{\lambda^{\alpha+2}},
\]
we get
\[
\mathbb{E}(X^2)
=
\frac{\lambda^\alpha}{\Gamma(\alpha)}
\cdot
\frac{\Gamma(\alpha+2)}{\lambda^{\alpha+2}}.
\]

Since
\[
\Gamma(\alpha+2)
=
(\alpha+1)\alpha\Gamma(\alpha),
\]
we have
\[
\mathbb{E}(X^2)
=
\frac{\alpha(\alpha+1)}{\lambda^2}.
\]

Therefore,
\[
\operatorname{Var}(X)
=
\mathbb{E}(X^2)-(\mathbb{E}(X))^2
=
\frac{\alpha(\alpha+1)}{\lambda^2}
-
\frac{\alpha^2}{\lambda^2}.
\]

Hence
\[
\operatorname{Var}(X)=\frac{\alpha}{\lambda^2}.
\]
\end{frame}





\begin{frame}{Gamma as a Sum of Exponentials}
Suppose
\[
X_1,\dots,X_n
\]
are independent Exponential$(\lambda)$ random variables.

Define
\[
T=X_1+\cdots+X_n.
\]

Then:
\[
{
T\sim \operatorname{Gamma}(n,\lambda).
}
\]




\begin{keybox}{Interpretation}
Gamma distributions describe waiting times for multiple arrivals.
\end{keybox}
\end{frame}



\begin{frame}{Gamma and Negative Binomial}
There is a strong analogy:

\vspace{0.25cm}

\begin{center}
\renewcommand{\arraystretch}{1.4}
\begin{tabular}{ccc}
\toprule
Discrete & Continuous & Interpretation \\
\midrule
Geometric & Exponential & waiting for first event \\
Negative Binomial & Gamma & waiting for multiple events \\
\bottomrule
\end{tabular}
\end{center}

\vspace{0.25cm}

\begin{keybox}{Big picture}
Gamma distributions are continuous analogs of Negative Binomial distributions.
\end{keybox}
\end{frame}
\begin{frame}{MGF of Gamma Distribution}
If
\[
X\sim \operatorname{Gamma}(\alpha,\lambda),
\]
then its MGF is
\[
M_X(t)
=
\left(
\frac{\lambda}{\lambda-t}
\right)^\alpha,
\]
for \(t<\lambda\).
\end{frame}



\begin{frame}{Proof: MGF of Gamma Distribution}
The density is
\[
f_X(x)
=
\frac{\lambda^\alpha}{\Gamma(\alpha)}
x^{\alpha-1}e^{-\lambda x},
\qquad x>0.
\]

By definition,
\[
M_X(t)
=
\mathbb{E}(e^{tX})
=
\int_0^\infty e^{tx} f_X(x)\,dx.
\]

Therefore,
\[
M_X(t)
=
\frac{\lambda^\alpha}{\Gamma(\alpha)}
\int_0^\infty
x^{\alpha-1}e^{tx}e^{-\lambda x}\,dx.
\]

Combining the exponentials, we have
\[
M_X(t)
=
\frac{\lambda^\alpha}{\Gamma(\alpha)}
\int_0^\infty
x^{\alpha-1}e^{-(\lambda-t)x}\,dx.
\]

Using the Gamma integral,
\[
\int_0^\infty x^{\alpha-1}e^{-cx}\,dx
=
\frac{\Gamma(\alpha)}{c^\alpha},
\qquad c>0,
\]
with \(c=\lambda-t\), we get
\[
M_X(t)
=
\left(
\frac{\lambda}{\lambda-t}
\right)^\alpha,
\qquad \text{for } \lambda-t>0.
\]
\end{frame}
\begin{frame}{Relationship Between Beta and Gamma}
If
\[
X\sim \operatorname{Gamma}(a,1),
\qquad
Y\sim \operatorname{Gamma}(b,1),
\]
independently, then
\[
{
\frac{X}{X+Y}
\sim
\operatorname{Beta}(a,b).
}
\]

\vspace{0.25cm}

\begin{keybox}{Beautiful connection}
Beta and Gamma distributions are deeply linked.
\end{keybox}
\end{frame}

\begin{frame}{Summary}
\begin{itemize}
    \item Beta distributions model random quantities between
    \[
    0
    \quad \text{and} \quad
    1.
    \]

    \item Gamma functions generalize factorials.

    \item Gamma distributions generalize Exponential distributions.

    \item Gamma distributions arise from sums of independent Exponentials.

    \item Beta and Gamma distributions are closely connected.
\end{itemize}

\vspace{0.25cm}

\begin{keybox}{Next lecture}
Conditional expectation and prediction.
\end{keybox}
\end{frame}